\chapter{Introduction}

\section{Background of the Study}
A defining characteristic of the global counselling landscape is the rapid transition to Digital Health Interventions (DHIs) for delivering psychological programs \citep{Wind2020_98, Naslund2017_76, Sutherland2018_77}. Traditional in-person counselling often presents insurmountable barriers for caregivers of children with Autism Spectrum Disorder (ASD), including high financial costs, severe time constraints, and logistical challenges. To democratize access to mental health support, international counselling programs increasingly utilize asynchronous e-learning, real-time videoconferencing, and mobile health (mHealth) applications to deliver targeted therapy. 

In Malaysia, the drive toward targeted counselling through specialized frameworks is not merely a clinical preference but a systemic necessity. The prevalence of ASD in Malaysia mirrors the global surge \citep{Santomauro2025_73}. However, the Malaysian experience is filtered through a complex lens of cultural stigma, varying socio-economic status, and a centralized healthcare system that often struggles with "one-size-fits-all" support models. For many Malaysian caregivers, spiritual coping and communal support (the "Kampung Spirit") are as vital as clinical psychotherapy. A critical analysis suggests that targeted counselling in Malaysia is only effective if data-driven groupings include these specific cultural belief systems as clustering indicators. 

\section{Problem Statement}
With a global pediatric prevalence of 0.77\%, Malaysian researchers face a persistent challenge with under-reporting in rural regions. Critically, the lack of a mandatory national registry for neurodivergence creates a "data shadow," where many families remain outside the reach of data-driven interventions. This creates a bifurcation in the caregiver population: urban families with access to private, data-rich interventions, and rural families relying on generalized, often outdated public support. 

Furthermore, the "Oxygen Mask Principle" faces a harsh reality in the Malaysian public sector. With a specialist-to-patient ratio that remains below recommended international levels, the caregiver burden is exacerbated by long wait times for diagnosis and intervention. When a parent's psychological health is compromised, it hinders the child's trajectory \citep{Rusu2025_94, SnchezAmate2024_95, Palmer2020_83}. In Malaysia, this creates a vicious cycle where the lack of targeted support for the parent leads to developmental plateauing in the child, further increasing the caregiving load. Ignoring the heterogeneity of these caregivers and forcing them into standardized counseling programs invariably leads to high attrition rates \citep{Wang2025_74, Liu2025_72}.

\section{Research Objectives}
\begin{enumerate}
    \item \textbf{To identify} the dominant psychometric and logistical characteristics of Malaysian ASD caregivers by estimating latent variable scores derived from the Theory of Planned Behavior (TPB).
    \item \textbf{To evaluate and compare} the statistical efficacy of three distinct unsupervised clustering algorithms (K-Means, Latent Profile Analysis, and Two-Stage Clustering) in segmenting the caregiver profiles.
    \item \textbf{To develop} a Targeted Virtual Counseling Model based on the specific qualitative needs of the mathematically optimal caregiver clusters.
\end{enumerate}

\section{Research Questions}
\begin{enumerate}
    \item What are the dominant psychometric and logistical characteristics influencing Malaysian ASD caregivers' engagement with support systems?
    \item Which unsupervised clustering algorithm (K-Means, Latent Profile Analysis, or Two-Stage Clustering) provides the mathematically optimal fit for profiling ASD caregivers?
    \item What does a Targeted Virtual Counseling Model look like when matched to the specific psychometric needs of the algorithmically discovered clusters?
\end{enumerate}

\section{Scope of the Study}
This thesis is specifically scoped to the computational and statistical analysis of caregiver psychology. It utilizes a secondary dataset collected as part of a broader, collaborative multi-disciplinary research initiative. The primary data collection (including ethical clearances and survey distribution) was conducted under the auspices of the collaborative project. Therefore, the methodological boundary of this study formally begins at the data pre-processing phase, focusing strictly on the application of Partial Least Squares Structural Equation Modeling (PLS-SEM) and unsupervised machine learning algorithms to the provided dataset.

\section{Significance of the Study}
The transition to Digital Health Interventions (DHIs) has been particularly transformative. The democratization of access to mental health equity can be realized through localized platforms that provide targeted counselling. This research provides significant clinical, policy, and theoretical contributions by establishing a mathematically rigorous framework for this democratization. By allowing for data-driven grouping at scale, clinicians in urban centers can provide precision support to parents in rural or remote areas. Ultimately, it transitions support from a generalized approach to a targeted framework, reducing caregiver burnout and bridging the mental health digital divide in Malaysia.
